\documentclass[frames,pdf,slideColor,colorBG,total]{prosper}
%\documentclass[frames,pdf,slideColor,colorBG,accumulates,total]{prosper}
\usepackage[francais]{babel}
\usepackage{amsmath, amsbsy, pstricks, pst-node, pst-text, pst-3d}
%\usepackage{amsmath, amsfont, amsbsy, pstricks, pst-node, pst-text, pst-3d}
\usepackage{moreverb,epsfig,color,subfigure}
\NoFrenchBabelItemize
\usepackage{ucs}
\usepackage[utf8x]{inputenc}
%\usepackage{graphicx,color,caption2,amssymb,pstricks,lmodern}

\newcommand{\items}[1]{\begin{itemize} \item #1 \end{itemize}}
\newcommand{\itemss}[2]{\begin{itemize} \item #1 \item #2 \end{itemize}}
\newcommand{\itemsss}[3]{\begin{itemize} \item #1 \item #2 \item #3 \end{itemize}}
\newcommand{\itemssss}[4]{\begin{itemize} \item #1 \item #2 \item #3 \item #4 \end{itemize}}
\newcommand{\slidetextsize}{\footnotesize}
%\newcommand{\captionstyle}[1]{\scriptsize \textbf{#1}}

\myitem{1}{\includegraphics[width=.4cm]{red-bullet-on-white}}
\myitem{2}{\includegraphics[width=.4cm]{green-bullet-on-white}}
\myitem{3}{\includegraphics[width=.4cm]{yellow-bullet-on-blue}}

% ================================
% End of Commands Used in Document
% ================================
\newlength{\MiniPageLeft}
\newlength{\MiniPageRight}


\usepackage{relsize}
\usepackage{lipsum}

%c from texinfo.tex
\def\ifmonospace{\ifdim\fontdimen3\font=0pt }

%c C plus plus
\def\C++{%
\ifmonospace%
    C++%
\else%
    C\kern-.1667em\raise.30ex\hbox{\smaller{++}}%
\fi%
\spacefactor1000 }


\title{Project Management}
\subtitle{Réponse à l'appel d'offre \textit{SLIPWAY}}
%\caption{IR}
\author{B.\textsc{Bizet} A.\textsc{Buisson} A.\textsc{Marguerite} R.\textsc{Rincé}}

%\Logo{\includegraphics[scale=0.40]{img/logouniv.eps}}
\Logo(9,-0.70){\includegraphics[scale=0.4]{img/logoPSI.eps}}

\institution{
  PSI entreprise \\
  4 rue de la Houssinière, \\
  BP92208, F-44322 Nantes cedex 03, FRANCE
}


\begin{document}
% make the title slide
\maketitle

%%%%%%%%%%%%%%%%%%%%%%%%%%%%%%%%%%%%%%%%%%%%%%%%%%%%%%%%%%%%%%%%%%%%%%%%%%%%
%\overlays{5}{
  \begin{slide}[Box]{Déroulement}
%    \textbf{}
%    \vspace{0,8cm}
      \begin{enumerate}
        \item Présentation
        \item Planning détaillée
        \item Plan de charges
        \item Journal d'actions
        \item Facture et plan de facturation
        \item Conclusion
      \end{enumerate}
  \end{slide}
%}
%%%%%%%%%%%%%%%%%%%%%%%%%%%%%%%%%%%%%%%%%%%%%%%%%%%%%%%%%%%%%%%%%%%%%%%%%%%%
\overlays{4}{
  \begin{slide}[Box]{Présentation de \textsc{PSI}}
    \textbf{Power Services Informatique}
    \vspace{0,3cm}
      \begin{Itemize}
        \fromSlide{2}{\item PSI est un société basée à Nantes. Fondée en 2009 par Anatole \textsc{Bizet}.}
        \fromSlide{3}{\item Dirigée aujourd'hui par JP.\textsc{Bizet}.}
        \fromSlide{4}{\item Jeune entreprise dynamique composée de 7 salariés, avec un chiffre d'affaire de 850 000 euros.}
        \fromSlide{5}{\item Ses collaborateurs : SFR, IBM, Total\dots}
      \end{Itemize}
  \end{slide}
}




%%%%%%%%%%%%%%%%%%%%%%%%%%%%%%%%%%%%%%%%%%%%%%%%%%%%%%%%%%%%%%%%%%%%%%%%%%%
\overlays{5}{
  \begin{slide}[Box]{Présentation de l'équipe}
    \textbf{L'équipe \textit{PSI} pour ce projet}
    \vspace{0,3cm}
    \begin{Itemize}
      \fromSlide{2}{\item Chef de projet JP.\textsc{Bizet} : Ancien de chef de projet chez Accenture (15 ans).}
      \fromSlide{3}{\item Ingénieur JM.\textsc{Lafrange} : Ancien expert C/\C++/python chez Logica (7 ans).}
      \fromSlide{4}{\item Technicien T.\textsc{Dupont} : Diplomé DUT GEII (Lyon).}
      \fromSlide{5}{\item Technicien D.\textsc{Baudouin} : BTS Informatique (Nantes).}
    \end{Itemize}
   \end{slide}
}

 %%%%%%%%%%%%%%%%%%%%%%%%%%%%%%%%%%%%%%%%%%%%%%%%%%%%%%%%%%%%%%%%%%%%%%%%%%%% 

\overlays{4}{
  \begin{slide}[Box]{Conclusion}
    \textbf{Pourquoi \textit{PSI} ?}
    \vspace{0,3cm}
      \begin{Itemize}
        \fromSlide{2}{\item Une entreprise à l'échelle du projet.}
% 
        \fromSlide{3}{\item Une entreprise souple et dynamique.}
 %       \vspace{0,3cm}
        \fromSlide{4}{\item Une entreprise expérimentée.}
      \end{Itemize}
  \end{slide}
}
%%%%%%%%%%%%%%%%%%%%%%%%%%%%%%%%%%%%%%%%%%%%%%%%%%%%%%%%%%%%%%%%%%%%%%%%%%%%%%ù%

%% \begin{slide}{Realpaver : exemple de sortie}
%% \hypertarget{code}{ }
%% \tiny\vspace{-.35in}
%% \begin{verbatim}
%% RealPaver v. 0.4 (c) LINA 2004
%% INITIAL BOX
%% x in ]-oo , +oo[
%% y in ]-oo , +oo[
%% OUTER BOX 1
%% x in [0.7924334198270921 , 0.795689483022687]
%% y in [0.8806243697296342 , 0.8872330220900007]
%% ...
%% OUTER BOX 3053
%% x in [-0.7956894830226868 , -0.7924334198270919]
%% y in [-0.8872330220900011 , -0.8806243697296345]
%% precision: 0.00661, elapsed time: 150 ms
%% END OF SOLVING
%% Property: reliable process (no solution is lost)
%% Elapsed time: 150 ms
%% \end{verbatim}
%% \end{slide}
%%  %%%%%%%%%%%%%%%%%%%%%%%%%%%%%%%%%%%%%%%%%%%%%%%%%%%%%%%%%%%%%%%%%%%%%%%%%%%% 
%% \overlays{4}{
%%   \begin{slide}[Box]{Définition d'une boîte}
%%        \slidetextsize
%%     \textbf{Entité atomique du pavage}
%%     \begin{Itemize}
%%       \fromSlide{2}{\item Un identifiant 
%%         \begin{Itemize}
%%         \item NUMBER
%%         \item STRING
%%         \end{Itemize}}
%%       \fromSlide{3}{\item Une liste de coordonnées
%%         \begin{Itemize}
%%         \item INTERVAL
%%       \end{Itemize}}
%%       \fromSlide{4}{\item Une liste des caractéristiques
%%         \begin{Itemize}
%%         \item NUMBER
%%         \item STRING
%%         \item INTERVAL
%%       \end{Itemize}}
%%     \end{Itemize}
%%   \end{slide}
%% }
%% %% %%%%%%%%%%%%%%%%%%%%%%%%%%%%%%%%%%%%%%%%%%%%%%%%%%%%%%%%%%%%%%%%%%%%%%%%%%%%
%% \overlays{4}{
%%   \begin{slide}[Box]{Étude des structures de données pour une boîte}
%%     \textbf{Structure de données}
%%     \begin{Itemize}
%%       \fromSlide{2}{\item Un tableau}
%%       \fromSlide{3}{\item Une liste}
%%       \fromSlide{4}{\item Une table de hachage}
%%     \end{Itemize}
%%   \end{slide}
%% }
%% %%%%%%%%%%%%%%%%%%%%%%%%%%%%%%%%%%%%%%%%%%%%%%%%%%%%%%%%%%%%%%%%%%%%%%%%%%%%
%%   \begin{slide}[Box]{Étude des structures de données pour une boîte}
%%     \textbf{Tests expérimentaux}
%%     \begin{table}[h]
%%       \centering
%%       \begin{tabular}{|c|c|c|c|}
%%         \hline
%%         \backslashbox{Structure} {Nombre d'accès} & $10^6$ & $10^7$ & $10^8$ \\
%%         \hline
%%        3 HashMap & 0.31s& 2.86s & 28.56s\\
%%        \hline
%%        3 ArrayList & 0.20s & 1.81s & 17.49s\\
%%         \hline
%%        1 ArrayList & 0.19s & 1.58s & 15.12s\\
%%         \hline
%%        1 LinkedList & 0.31s & 2,79s &  27.59s\\
%%         \hline
%%       \end{tabular}
%%     \end{table}
%%     Relevé des temps CPU d'accès au caractéristiques en secondes
%%   \end{slide}
%% %%%%%%%%%%%%%%%%%%%%%%%%%%%%%%%%%%%%%%%%%%%%%%%%%%%%%%%%%%%%%%%%%%%%%%%%%%%%
%% \overlays{5}{
%%   \begin{slide}[Box]{Étude des structures de données pour un pavage}
%%     \textbf{Structure de données}
%%     \begin{Itemize}
%%       \fromSlide{2}{\item Vector}
%%       \fromSlide{3}{\item Dictionnaire}
%%       \fromSlide{4}{\item Arbre binaire}
%%       \fromSlide{5}{\item Arbre a-b}
%%     \end{Itemize}
%%   \end{slide}
%% }
%% %%%%%%%%%%%%%%%%%%%%%%%%%%%%%%%%%%%%%%%%%%%%%%%%%%%%%%%%%%%%%%%%%%%%%%%%%%%%
%%   \begin{slide}[Box]{Visualisation : QuadTree}
%%     \textbf{Exemple de représentation}
%%       \begin{figure}[t]
%%         \begin{center}
%%           \includegraphics[width=0.55\textwidth]{img/quadtree}
%%         \end{center}
%%       \end{figure}
%%   \end{slide}
%% %% %%%%%%%%%%%%%%%%%%%%%%%%%%%%%%%%%%%%%%%%%%%%%%%%%%%%%%%%%%%%%%%%%%%%%%%%%%%%
%%   \begin{slide}[Box]{Visualisation : QuadTree}
%%     %    \slidetextsiz
%%     \textbf{Algorithme}
%%     \begin{table}[htbp]
%%       \centering
%%       \begin{tabular}{|c|cccc|}
%%         \hline
%%    Cas de récursion & \includegraphics[scale=0.20]{img/QT4}&  \includegraphics[scale=0.20]{img/QT5}& &\\
%%    \hline
%%    Cas d'arrêt&  \includegraphics[scale=0.20]{img/QT1}&  \includegraphics[scale=0.20]{img/QT2}&  \includegraphics[scale=0.20]{img/QT3}&  \includegraphics[scale=0.30]{img/QT6}\\
%%    \hline
%%       \end{tabular}
%%     \end{table}
%% Classification visuelle des cas d'arrêts et de récursions

%%   \end{slide}
%%   %% %%%%%%%%%%%%%%%%%%%%%%%%%%%%%%%%%%%%%%%%%%%%%%%%%%%%%%%%%%%%%%%%%%%%%%%%%%%%
%% \overlays{3}{
%%   \begin{slide}[Box]{Visualisation : QuadTree}
%%     %    \slidetextsiz
%%     \textbf{Inconvénients du QuadTree}
%%     \begin{Itemize}
%%       \fromSlide{2}{\item Structure peu adaptée pour stocker des éléments de dimension $d$}
%%       \fromSlide{3}{\item Problèmes des \emph{frontières}}
%%     \end{Itemize}
%%   \end{slide}
%% }
%%   %%%%%%%%%%%%%%%%%%%%%%%%%%%%%%%%%%%%%%%%%%%%%%%%%%%%%%%%%%%%%%%%%%%%%%%%%%%%

%%   \begin{slide}[Box]{Visualisation : QuadTree}
%%       \textbf{Cas de superposition des boîtes}
%%       \begin{figure}[t]
%%         \begin{center}
%%           \includegraphics[width=0.25\textwidth]{img/QT8}
%%         \end{center}
%%       \end{figure}
%%   \end{slide}
%%   %%%%%%%%%%%%%%%%%%%%%%%%%%%%%%%%%%%%%%%%%%%%%%%%%%%%%%%%%%%%%%%%%%%%%%%%%%%%

%% \overlays{6}{
%%   \begin{slide}[Box]{Visualisation : QuadTree}
%%       \textbf{Le problème des \emph{frontières}}
%%    \onlySlide*{1}{
%%    \begin{figure}[t]
%%         \begin{center}
%%           \includegraphics[width=0.50\textwidth]{img/SQT1}
%%         \end{center}
%%       \end{figure}
%% }
%%    \onlySlide*{2}{
%%    \begin{figure}[t]
%%         \begin{center}
%%           \includegraphics[width=0.50\textwidth]{img/SQT2}
%%         \end{center}
%%       \end{figure}
%% }
%%    \onlySlide*{3}{
%%      \begin{figure}[t]
%%        \begin{center}
%%          \includegraphics[width=0.50\textwidth]{img/SQT3}
%%        \end{center}
%%      \end{figure}
%%    }
%%    \onlySlide*{4}{
%%    \begin{figure}[t]
%%         \begin{center}
%%           \includegraphics[width=0.50\textwidth]{img/SQT4}
%%         \end{center}
%%       \end{figure}
%% }
%%    \onlySlide*{5}{
%%    \begin{figure}[t]
%%         \begin{center}
%%           \includegraphics[width=0.50\textwidth]{img/SQT5}
%%         \end{center}
%%       \end{figure}
%% }
%%    \fromSlide{6}{
%%    \begin{figure}[t]
%%         \begin{center}
%%           \includegraphics[width=0.50\textwidth]{img/SQT6}
%%         \end{center}
%%       \end{figure}
%% }

%%  \end{slide}
%% }
%%   %%%%%%%%%%%%%%%%%%%%%%%%%%%%%%%%%%%%%%%%%%%%%%%%%%%%%%%%%%%%%%%%%%%%%%%%%%%%
%%   \begin{slide}[Box]{Visualisation : R-Tree}
%%           \textbf{Exemple de représentation}
%%       \begin{figure}[t]

%%         \begin{center}
%%           \includegraphics[width=0.55\textwidth]{img/rtree}
%%         \end{center}
%%       \end{figure}
%%   \end{slide}
%%   %%%%%%%%%%%%%%%%%%%%%%%%%%%%%%%%%%%%%%%%%%%%%%%%%%%%%%%%%%%%%%%%%%%%%%%%%%%%
%% \overlays{4}{
%%   \begin{slide}[Box]{Visualisation : R-Tree}
%%     \textbf{Description du R-Tree}
%%     \begin{Itemize}
%%       \fromSlide{2}{   \item Variante de l'arbre B}
%%       \fromSlide{3}{    \item Chaque nœud est un MBR de ces fils}
%%       \fromSlide{4}{  \item Chaque nœud contient entre $m$ et $M$ fils \\ avec $m \leq \frac{M}{2}$}
%%     \end{Itemize}
%%   \end{slide}
%% }
%%   %%%%%%%%%%%%%%%%%%%%%%%%%%%%%%%%%%%%%%%%%%%%%%%%%%%%%%%%%%%%%%%%%%%%%%%%%%%%

%%   \begin{slide}[Box]{Visualisation : R-Tree}
%%       \textbf{Hauteur de l'arbre et nombre de nœuds au pire cas}
%%     \begin{Itemize}
%%     \item Hauteur
%%       \begin{equation*}
%%         h_{max} = \left\lceil\log_{m} n\right\rceil - 1
%%       \end{equation*}
%%       \item Nombre de nœuds
%%       \begin{equation*}
%%         N_{noeuds}=\sum_{i=1}^{h_{max}}{\left\lceil\frac{n}{m^i}\right\rceil}
%%       \end{equation*}
%%     \end{Itemize}
%%   \end{slide}

%%   %%%%%%%%%%%%%%%%%%%%%%%%%%%%%%%%%%%%%%%%%%%%%%%%%%%%%%%%%%%%%%%%%%%%%%%%%%%%

%% \begin{slide}[Box]{Que reste-t-il à faire ?}
%%     \vspace{1cm}
%%     \begin{Itemize}
%%     \item Chercher la meilleure implémentation pour le R-Tree.
%%     \item Une étude approfondie pour prendre en compte la gestion des filtres.
%%     \item Trouver un moyen efficace de sauvegarde.
%%     \end{Itemize}
%%   \end{slide}

\end{document}
